\documentstyle[%


righttag%


,11pt%


,amssymb%


,russian%               remove in Eng.


,izvru%                 Set izven% in Eng.


]{amsart}





% NB for eng.: replace separating , for .





\newcommand{\Grad}{\operatorname{grad}}


\newcommand{\Div}{\operatorname{div}}


\newcommand{\Ker}{\operatorname{Ker}}


\newcommand{\im}{\operatorname{Im}}


\newcommand{\tts}[1]{}





\theoremstyle{plain}


\theorembodyfont{\it}


\newtheorem{lemnonum}{\hskip\parindent Лемма}


\renewcommand{\thelemnonum}{}





%\hoffset=-15mm%                  For Eng.


 


\setcounter{page}{3}


\god{1999}


\nomer{5~(444)} %                        Remove in Eng.


%\nomer{Vol.\,43, No.\,5}%               Set in Eng.


%\str{\pageref{begin}--\pageref{end}}%  Set in Eng.


\udk{519.63}





\author{\it В.Н.\,Абрашин, В.М.\,Волков, А.А.\,Егоров, Н.Г.\,Жадаева}


% NB:   ^^ remove in Eng.


\title{Об одном классе разностных методов решения\\


уравнений Навье--Стокса}





\date{}


%\pagestyle{myheadings}%         Set in Eng.


\pagestyle{plain} %              Remove in Eng.


\begin{document}


%\thispagestyle{plain}%          Set in Eng.


\maketitle


\label{begin}





Работа посвящена вопросам построения разностных схем и итерационных


методов для решения $(u,v,P)$-систем вязкой несжимаемой жидкости.


Основное внимание уделяется решению линеаризованных уравнений


Навье--Стокса с использованием идеи многокомпонентного метода


переменных направлений (ММПН) [1]--[4], примыкающего к методам расщепления


[5], [7] и обладающего полной аппроксимацией [8], [9]. В качестве базовых


разностных схем используются энергетически нейтральные схемы,


предложенные


для уравнений Навье--Стокса в [10]. Выбор схем такого вида объясняется


возможностью производить вычисления без постановки дополнительных


условий


для давления на границе или в нерасчетной области (что, вообще говоря,


естественно) [11], [12] с точным удовлетворением условиям прилипания на


границе и разностному уравнению неразрывности. При построении


экономичных


методов сохранить эти условия не всегда удается. Это касается как метода


расщепления [6], [7], так и метода переменных направлений [13]. В связи


с этим следует отметить работы [14], [15], в которых на разнесенных сетках


рассматривается явная схема расщепления по физическим факторам,


позволяющая


удовлетворить условия прилипания с высокой точностью.





Энергетически нейтральные разностные схемы являются неявными алгоритмами


довольно сложной конструкции, поэтому для их реализации требуется


построение экономичных методов, которые сохраняли бы их свойства


консервативности и полной аппроксимации.





В данной работе предложены экономичные методы полной аппроксимации для


расщепления уравнений по физическим процессам и по размерности,


а также методы декомпозиции области [16]. Изучены как безитерационные


методы для нестационарных задач, так и итерационные процессы для решения


стационарных и нестационарных уравнений динамики вязкой несжимаемой


жидкости. Большое внимание уделяется вопросам распараллеливания


вычислительного алгоритма [17], [18].





В области $G_T=\overline G\times [0,T]$, $\overline G=


\{x=(x_1,x_2),\ 0\leq x_\alpha\leq l_\alpha,\ \alpha=1,2\}$ ---


прямоугольник с границей $\Gamma$, рассмотрим


линеаризованную двумерную краевую задачу для уравнений Навье--Стокса


в переменных скорость--давление


\begin{gather}


\frac{\partial u_k}{\partial t}-\nu\sum_{\alpha=1}^2\frac{\partial^2u_k}{


\partial x_\alpha^2}+\sum_{\alpha=1}^2a_\alpha\,\frac{\partial


u_k}{\partial


x_\alpha}=-\frac{\partial P}{\partial x_k},\quad


(x,t)\in G\times (0,T],\quad


\nu>0,\quad k=1,2,\\


\sum_{\alpha=1}^2\frac{\partial u_\alpha}{\partial x_\alpha}=0,\quad


(x,t)\in G\times (0,T],\\


u_k(x,0)=\varphi_k(x),\quad x\in \overline G,\quad


u_k(x,t)=0,\quad (x,t)\in \Gamma\times (0,T].


\end{gather}


В уравнении (1) коэффициенты $a_\alpha$, $\alpha=1,2$, --- некоторые


ограниченные константы. В данной постановке из [19] следует


существование


и единственность решения задачи (1)--(3) (давление $P$ определяется


с точностью до аддитивной постоянной).





Для построения численного алгоритма в области $G_T$ введем две


равномерные


по каждому направлению пространственные сетки: сетку целых узлов


$\overline\omega_h=\{x=(x_1^{(i_1)},x_2^{(i_2)})=(i_1h_1,i_2h_2),$


$i_\alpha=0,\dotsc,N_\alpha,\ h_\alpha=l_\alpha/N_\alpha,\ \alpha=1,2\}$


с границей


$\gamma_h=\overline\omega_h\setminus\omega_h$ и сетку полуцелых узлов


$\Omega_h=\{x=(x_1^{(i_1-0,5)},x_2^{(i_2-0,5)})=((i_1-0,5)h_1,(i_2-0,5)h_2),\ 


i_\alpha=1,\dotsc,N_\alpha,\ h_\alpha=l_\alpha/N_\alpha,\ \alpha=1,2\}$;


временную сетку определим стандартным образом $\overline\omega_\tau=


\{t_n=n\tau,\ n=0,1,\dotsc\}$. Кроме того, разобьем $\Omega_h$ на два


сеточных подмножества: $\Omega_h=S_h\cup T_h$,


$S_h=\{x=(x_1^{(i_1-0,5)},x_2^{(i_2-0,5)}),\


i_\alpha=2,\dotsc,N_\alpha-1,\ \alpha=1,2\}$,


$T_h=\{x=(x_1^{(i_1-0,5)},x_2^{(i_2-0,5)}),\ i_1=1,N_1,\


i_2=1,\dotsc,N_2;\


i_2=1,N_2,\ i_1=2,\dotsc,N_2-1\}$.





Пусть $H_1$ --- пространство сеточных функций, определенных на


$\overline\omega_h$ и обращающихся в нуль на $\gamma_h$, со скалярным


произведением


$(u,v)_{H_1}=\sum\limits_{x\in\omega_h}h_1h_2u(x)v(x)$ и


нормой


$\|u\|_{1}=(u,u)_{H_1}^{1/2}$. Аналогичным образом определим


пространство


$H_2$ сеточных функций, заданных на сетке $\Omega_h$, со скалярным


произведением


$[u,v]_{H_2}=\sum\limits_{x\in\Omega_h}\ov h_1\ov


h_2u(x)v(x)$


и нормой $\|u\|_{2}=[u,u]_{H_2}^{1/2}$, где


$\ov h_\alpha=\begin{cases}


h_\alpha, & i_\alpha=2,\dotsc,N_\alpha-1; \\


0,5h_\alpha, & i_\alpha=1,N_\alpha,\quad \alpha=1,2.\end{cases}$


Далее в $H_1$ будем рассматривать сеточные аналоги скоростей $u_1$, $u_2$,


а в $H_2$ --- сеточные аналоги давления $P$.





Помимо общепринятых обозначений теории разностных схем будем


использовать


также следующие обозначения (для упрощения записи вместо $i_1$, $i_2$ стоит


$i$,~$j$):


\begin{gather*}


u_{\wh x_1}^{i,j}=(u_{i+1/2,j+1/2}+u_{i+1/2,j-1/2}-


u_{i-1/2,j+1/2}-u_{i-1/2,j-1/2})/(2h_1),\quad


u\in\Omega_h,\\


u_{\wh x_2}^{i,j}=(u_{i+1/2,j+1/2}+u_{i-1/2,j+1/2}-


u_{i+1/2,j-1/2}-u_{i-1/2,j-1/2})/(2h_2),\quad u\in\Omega_h, \\


v_{\wt x_1}^{i-1/2,j-1/2}=(v_{i,j}+v_{i,j-1}-


v_{i-1,j}-v_{i-1,j-1})/(2h_1),\quad v\in\omega_h,\\


v_{\wt x_2}^{i-1/2,j-1/2}=(v_{i,j}+v_{i-1,j}-


v_{i,j-1}-v_{i-1,j-1})/(2h_2),\quad v\in\omega_h.


\end{gather*}





Заменяя функции $u_1$, $u_2$, $P$


их сеточными аналогами $y_1$, $y_2\approx u_1$,


$z_1$, $z_2\approx u_2$, $q\approx P$, аппроксимируем задачу (1)--(3)


следующей разностной схемой ММПН:


\begin{equation}


\begin{array}{c}


Y_{1t}=-\Grad_h\widehat Q+LY_2,\quad


(x,t)\in\omega_h\times\omega_\tau,\\


Y_{2t}=-\Grad_h\widehat Q+L\widehat Y_2,\quad


(x,t)\in\omega_h\times\omega_\tau, \\


\Div_h\widehat Y_1\equiv\wh y_{1\wt x_1}+\wh z_{1\wt x_2}=0,


\quad (x,t)\in \Omega_h\times\omega_\tau,\\


y_\alpha(x,0)=\varphi_1(x),\quad z_\alpha(x,0)=\varphi_2(x),\quad


x\in\overline\omega_h,\quad \alpha=1,2,\\


y_\alpha(x,t)=z_\alpha(x,t)=0,\quad


(x,t)\in\gamma_h\times\omega_\tau,\quad\alpha=1,2,


\end{array}


\end{equation}


где


$Y_{\alpha}=(y_\alpha,z_\alpha)^{*}$, $\alpha=1,2$,


$\Grad_hQ=(q_{\wh x_1},q_{\wh x_2})^*$,


$LY_2=(\Lambda y_2,\Lambda z_2)^*$, $\Lambda v=\nu(v_{\ov x_1x_1}+


v_{\ov x_2x_2})-a_1v_{\overset{\circ}{x}{}_1}-


a_2v_{\overset{\circ}{x}{}_2}$


(символ $*$ здесь обозначает транспонирование).





Алгоритм (4) реализуется в два этапа.


На первом этапе для определения давления $\wh q$


подставим выражения для компонент вектора


скорости


$\wh y_1,\wh z_1$


из первого уравнения (4) в уравнение неразрывности, после несложных


преобразований получим следующую задачу:


\begin{equation}


\wh q_{\wh x_1\wt x_1}+\wh q_{\wh x_2\wt x_2}=


F(y_2,z_2),\quad x\in S_h,\quad


\wh q_n=f(y_2,z_2),\quad x\in T_h,


\end{equation}


где


\begin{align*}


\wh q_n^{\,i-1/2,1/2}&=0,25(\wh q_{x_2}^{\,i+1/2,1/2}+


2\wh q_{x_2}^{\,i-1/2,1/2}+\wh q_{x_2}^{\,i-3/2,1/2})+


0,25h_2(\wh q_{\ov x_1x_1}^{\,i-1/2,3/2}+


\wh q_{\ov x_1x_1}^{\,i-1/2,1/2}),\\


%


\wh q_n^{\,i-1/2,N_2-1/2}&=


0,25(-\wh q_{\ov x_2}^{\,i+1/2,N_2-1/2}-


2\wh q_{\ov x_2}^{\,i-1/2,N_2-1/2}-\wh q_{\ov


x_2}^{\,i-3/2,N_2-1/2})+\\


&\qquad+0,25h_2(\wh q_{\ov x_1x_1}^{\,i-1/2,N_2-3/2}+


\wh q_{\ov x_1x_1}^{\,i-1/2,N_2-1/2}),\\


%


\wh q_n^{\,1/2,j-1/2}&=0,25(\wh q_{x_1}^{\,1/2,j+1/2}+


2\wh q_{x_1}^{\,1/2,j-1/2}+\wh q_{x_1}^{\,1/2,j-3/2})+


0,25h_1(\wh q_{\ov x_2x_2}^{\,3/2,j-1/2}+


\wh q_{\ov x_2x_2}^{\,1/2,j-1/2}),\\


%


\wh q_n^{\,N_1-1/2,j-1/2}&=


0,25(-\wh q_{\ov x_1}^{\,N_1-1/2,j+1/2}-


2\wh q_{\ov x_1}^{\,N_1-1/2,j-1/2}-\wh q_{\ov


x_1}^{\,N_1-1/2,j-3/2})+\\


&\qquad+0,25h_1(\wh q_{\ov x_2x_2}^{\,N_1-3/2,j-1/2}+


\wh q_{\ov x_2x_2}^{\,N_1-1/2,j-1/2}), \\


%


\displaybreak[3]


\wh q_n^{\,1/2,1/2}&=0,5h_2(\wh q_{x_1}^{\,1/2,3/2}+


\wh q_{x_1}^{\,1/2,1/2})/(h_1+h_2)+0,5h_1(\wh q_{x_2}^{\,3/2,1/2}+


\wh q_{x_2}^{\,1/2,1/2})/(h_1+h_2),\\


%


\wh q_n^{\,1/2,N_2-1/2}&=


0,5h_2(\wh q_{x_1}^{\,1/2,N_2-3/2}+


\wh q_{x_1}^{\,1/2,N_2-1/2})/(h_1+h_2)+\\


&\qquad+0,5h_1


(-\wh q_{\ov x_2}^{\,3/2,N_2-1/2}-


\wh q_{\ov x_2}^{\,1/2,N_2-1/2})/(h_1+h_2),\\


%


\wh q_n^{\,N_1-1/2,N_2-1/2}&=0,5h_2(-\wh q_{\ov


x_1}^{\,N_1-1/2,N_2-3/2}-


\wh q_{\ov x_1}^{\,N_1-1/2,N_2-1/2})/(h_1+h_2)+\\


&\qquad+0,5h_1


(-\wh q_{\ov x_2}^{\,N_1-3/2,N_2-1/2}-


\wh q_{\ov x_2}^{\,N_1-1/2,N_2-1/2})/(h_1+h_2),\\


%


\wh q_n^{\,N_1-1/2,1/2}&=0,5h_2(-\wh q_{\ov x_1}^{\,N_1-1/2,3/2}-


\wh q_{\ov x_1}^{\,N_1-1/2,1/2})/(h_1+h_2)+\\


&\qquad+0,5h_1(\wh q_{x_2}^{\,N_1-3/2,1/2}+


\wh q_{x_2}^{\,N_1-1/2,1/2})/(h_1+h_2),\\


%


&\qquad x\in T_h,\quad


F(y_2,z_2)=(\Lambda y_2)_{\wt x_1}


+(\Lambda z_2)_{\wt x_2},\quad x\in S_h,


\end{align*}


а функцию $f$ непосредственно получаем из уравнения неразрывности


в узлах $x\in T_h$ и первых двух уравнений движения в (4).





Таким образом, после первого этапа реализации алгоритма (4) имеем на


сетке


$\Omega_h$ значения $\wh q$ (с точностью до произвольного слагаемого),


которые, как будет показано ниже, определяют единственным образом на


$\omega_h$\; $\Grad_h\widehat Q$. Второе векторное уравнение


представляет


собой две пятиточечные разностные схемы для $\wh y_2$, $\wh z_2$,


принимаемые


в качестве приближенного решения к компонентам вектора скорости


$u_1,u_2$.


Эти схемы можно решать как прямыми, так и итерационными методами.





Изучим вопросы разрешимости задачи для давления. Рассмотрим разностную


схему


\begin{equation}


\Lambda y\equiv y_{\wh x_1\wt x_1}+


y_{\wh x_2\wt x_2}=F,\quad x\in S_h,\quad


y_n=f,\quad x\in T_h,


\end{equation}


являющуюся аналогом схемы (5). Введем следующие обозначения:


$U$ --- множество угловых точек $\Omega_h$;


$\Pi'$ ($\Pi''$) --- множество узлов из данного множества,


для которых $i+j$ четное (нечетное);


$S_h'$ ($S_h''$) --- множество четных (нечетных) узлов $S_h$;


$T_{1h}$ --- множество узлов левой и правой границ $T_h$,


$T_{2h}$ --- множество узлов нижней и верхней границ $T_h$.





\begin{lemnonum}


Для задачи $(6)$ имеют место равенства


\begin{align*}


\sum_{S_h}h_1h_2F+\sum_{T_{1h}\setminus U}h_2f+


\sum_{T_{2h}\setminus U}h_1f+\sum_U0,5(h_1+h_2)f&=0,\\


\sum_{S_h'}h_1h_2F+\sum_{T_{1h}'\setminus U'}h_2f+


\sum_{T_{2h}'\setminus U'}h_1f+\sum_{U'}0,5(h_1+h_2)f&=0,\\


\sum_{S_h''}h_1h_2F+\sum_{T_{1h}''\setminus U''}h_2f+


\sum_{T_{2h}''\setminus U''}h_1f+\sum_{U''}0,5(h_1+h_2)f&=0.


\end{align*}


\end{lemnonum}





Справедливость леммы проверяется непосредственной подстановкой в


рассматриваемые равенства вместо $F$ и $f$ их выражений из (5) и


является прямым следствием соотношений


$$


\sum_{\Omega_h}(u_{\wt x_1}+v_{\wt x_2})=0,\quad


\sum_{\Omega_h'}(u_{\wt x_1}+v_{\wt x_2})=0,\quad


\sum_{\Omega_h''}(u_{\wt x_1}+v_{\wt x_2})=0,\quad


u,\,v\in H_1.


$$





Далее, запишем задачу (6) в операторном виде


\begin{equation}


Ay=g,


\end{equation}


где $A=-\overline\Lambda$ --- линейный ограниченный самосопряженный


оператор в гильбертовом конечномерном пространстве $H_2$ [20].


Обратимся к теории операторных уравнений. Пусть


$\Ker A$ --- ядро оператора $A$, $\im A$ --- его образ.


Тогда пространство $H_2$ есть прямая сумма ортогональных подпространств:


$H_2=\Ker A+\im A^*$, $H_2=\Ker A^*+\im A$.


Известно, что уравнение (7) однозначно разрешимо при любой правой части


$g\in H_2$ тогда и только тогда, когда $\Ker A=0$. Если


$\Ker A\ne 0$, то для разрешимости неоднородного уравнения (7)


необходимо и достаточно, чтобы правая часть $g$ была ортогональна


подпространству $\Ker A^*$. В этом случае решение неединственно


и определяется с точностью до произвольного элемента, принадлежащего


$\Ker A$: $y=y_1+y_2$, $y_1\in\Ker A$, $Ay_2=g$,


$y_2\in\im A^*$. Построенное таким образом решение $y$ задачи (7)


называется классическим. Если, кроме того, $y$ имеет наименьшую норму


($\Ker A\ne 0$, $g$ ортогональна $\Ker A^*$), то такое


решение


единственно, принадлежит $\im A^*$ и называется нормальным.





Сформулируем теперь теорему существования классического решения


операторного уравнения (7), а значит, и разностной задачи (5),


возникающей


на первом этапе реализации алгоритма (4).





\begin{thm}


Лемма гарантирует существование в


классическом смысле решения задачи $(7)$.


\end{thm}





Теорема 1 доказывается непосредственной проверкой


с помощью равенств $[g,\mu_{11}]_{H_2}=0$ и


$[g,\mu_{N_1N_2}]_{H_2}=0$, где $\mu_{11},\mu_{N_1N_2}$


--- собственные функции оператора $A$ задачи (7),


образующие базис в подпространстве $\Ker A$


(все остальные собственные функции образуют базис в


подпространстве $\im A$).





Поскольку базис подпространства $\Ker A$ состоит из двух функций


$\mu_{11}$ и $\mu_{N_1N_2}$, а произвольный элемент, с точностью


до которого определяется решение задачи (7), принадлежит данному


подпространству, то этот элемент является некоторой линейной комбинацией


функций $\mu_{11},\mu_{N_1N_2}$. Он может распадаться на две


различные константы, одна из которых будет определена в узлах сетки


$\Omega_h'$, а другая --- в узлах сетки $\Omega_h''$. Однако при этом


градиенты $y_{\wh x_1},y_{\wh x_2}$, полученные на ячейках


сетки $\Omega_h$, будут определяться единственным образом. Единственное


нормальное решение $y\in\im A$ можно найти аналогично [20].





Таким образом, разностная схема (4) является примером применения


ММПН для расщепления по физическим процессам, при этом сохраняется


полная аппроксимация исходной задачи и выполняются классические условия


разрешимости.





Обозначим $\|u\|_{(1)}^2=\|u_{\ov x_1}]|_1^2+|[u_{x_1}\|_1^2


+\|u_{\ov x_2}]|_1^2+|[u_{x_2}\|_1^2$. Для устойчивости


разностной схемы (4) имеет место





\begin{thm}


Разностная схема $(4)$ безусловно устойчива по начальным данным,


и для ее решения справедлива оценка


$$


B(Y,Q)\leq C_2(\|\varphi_1\|_{(1)}^2+


\|\varphi_2\|_{(1)}^2+\|\Lambda\varphi_1\|^2+\|\Lambda\varphi_2\|^2),


$$


где


$B(Y,Q)=\|y_2\|_{(1)}^2+\|z_2\|_{(1)}^2+C_1(\|\Lambda y_2-


q_{\wh x_1}\|^2+\|\Lambda z_2-q_{\wh x_2}\|^2)$,


$C_1>0$, $C_2>0$ --- постоянные, не зависящие от сеточных шагов.


\end{thm}





При доказательстве теоремы~2 используется техника, применяемая


для получения аналогичных оценок устойчивости ММПН [1]--[4].


Так как разностная схема (4) удовлетворяет условию полной аппроксимации,


то


нетрудно показать, что при достаточной гладкости искомого решения


схема имеет скорость сходимости порядка $O(h^2+\tau)$. Ухудшение свойств


гладкости естественно приведет к уменьшению скорости сходимости метода.





Помимо разностной схемы (4) для решения задачи (1)--(3) можно также


использовать алгоритм асинхронного типа


\begin{equation}


\begin{array}{c}


Y_{1t}=-\sigma\Grad_h(\widehat Q-Q)-\Grad_hQ+LY_2, \\


Y_{2t}=\sigma L(\widehat Y_2-Y_2)-\Grad_hQ+LY_2,\quad


\Div_h\widehat Y_1=0.


\end{array}


\end{equation}


Для устойчивости разностной схемы (8) утверждение теоремы~2 сохраняется


при $\sigma\geq 1$. Очевидно, применение схемы (8) увеличивает


возможности распараллеливания вычислительного процесса.





На базе схемы (4) несложно построить экономичные


трехкомпонентные методы, например, вида


\begin{equation}


\begin{array}{c}


Y_{1t}=-\Grad_h\widehat Q+L_2Y_2+L_3Y_3,\quad


Y_{2t}=-\Grad_h\widehat Q+L_2\widehat Y_2+L_3Y_3,\\


Y_{3t}=-\Grad_h\widehat Q+L_2\widehat Y_2+L_3\widehat Y_3,\quad


\Div_h\widehat Y_1=0,


\end{array}


\end{equation}


или


\begin{equation}


\begin{array}{c}


Y_{1t}=-\sigma\Grad_h(\widehat Q-Q)+M(Q,Y_2,Y_3),\\


Y_{\alpha t}=\sigma L_\alpha(\widehat


Y_\alpha-Y_\alpha)+M(Q,Y_2,Y_3),\quad


\alpha=2,3,\quad \Div_h\widehat Y_1=0,


\end{array}


\end{equation}


где


$M(Q,Y_2,Y_3)=-\Grad_hQ+L_2Y_2+L_3Y_3$,


$L_{\alpha+1}Y_{\alpha+1}=(\Lambda_{\alpha+1}y_{\alpha+1},


\Lambda_{\alpha+1}z_{\alpha+1})^*$,


$Y_{\alpha}=(y_\alpha,z_\alpha)^{*}$,


$\Lambda_{\alpha+1}y_{\alpha+1}=\nu y_{\alpha+1,\ov x_\alpha


x_\alpha}-a_\alpha y_{\alpha+1,\overset{\circ}{x}{}_\alpha}$,


$\Lambda_{\alpha+1}z_{\alpha+1}=\nu z_{\alpha+1,\ov x_\alpha


x_\alpha}-a_\alpha z_{\alpha+1,\overset{\circ}{x}{}_\alpha}$,


$\alpha=1,2$, $\Grad_hQ=(q_{\wh x_1},q_{\wh x_2})^*$.





Справедлива





\begin{thm}


Разностные схемы $(9)$, $(10)$ безусловно устойчивы по начальным данным


$($схема $(10)$ устойчива при $\sigma\geq 3/2)$, и для их решения выполняется


оценка


\begin{multline*}


\sum_{\alpha=1}^2(\|y_{\alpha\ov x_\alpha}]|_\alpha^2+


\|z_{\alpha\ov


x_\alpha}]|_\alpha^2)+C_3\bigg(\|y_{3t}\|^2+\|z_{3t}\|^2+


\bigg\|\sum_{\alpha=1}^2\Lambda_\alpha y_\alpha-q_{\wh


x_1}\bigg\|^2+


\bigg\|\sum_{\alpha=1}^2\Lambda_\alpha z_\alpha-q_{\wh


x_2}\bigg\|^2\bigg)\leq \\


\leq C_4\bigg(


\sum_{\alpha=1}^2(\|\varphi_{1\ov x_\alpha}]|_\alpha^2+


\|\varphi_{2\ov x_\alpha}]|_\alpha^2)+


\bigg\|\sum_{\alpha=1}^2\Lambda_\alpha\varphi_1-q_{0\wh


x_1}\bigg\|^2+


\bigg\|\sum_{\alpha=1}^2\Lambda_\alpha\varphi_2-q_{0\wh x_2}


\bigg\|^2\bigg),


\end{multline*}


где


$C_3,C_4=\mathrm{const}>0$.


\end{thm}





Методы (8), (10) имеют некоторые ограничения, связанные, например,


с количеством пространственных переменных. Для отыскания $Q$


используется,


как правило, прямой метод, основанный на быстром преобразовании Фурье.


Что касается задачи определения компоненты $Y_2$, то


дополнительные возможности в этом направлении имеет метод декомпозиции


пространственной области. Рассмотрим некоторые варианты этого метода


применительно к решению системы (8) при $\sigma=1$.





Разобьем область $\overline G$ на подобласти прямыми


$x_{1i}$, $i=1,\ldots,N$, параллельными оси $Ox_2$,


так, чтобы узлы сетки $x_{i_1}=i_1h_1$, $i_1=0,\ldots,N_1$, лежали


на соответствующих прямых. Тем самым сетка ${\overline\omega}_h$


разбивается на подобласти ${\overline\omega}_h^{\,i}$ (рис.~1).





\unitlength=1.00mm


\special{em:linewidth 0.4pt}


\linethickness{0.4pt}


\begin{picture}(20.00,75.00)


\put(40.00,40.00){\line(1,0){90.00}}


\put(130.00,40.00){\line(-1,0){90.00}}


\put(40.00,40.00){\line(0,1){30.00}}


\put(40.00,40.00){\line(1,0){90.00}}


\put(128.00,40.00){\vector(1,0){3.00}}


\put(40.00,68.00){\vector(0,1){3.00}}


\put(40.00,60.00){\line(1,0){80.72}}


\put(121.00,60.00){\line(0,-1){20.00}}


\put(55.00,60.00){\line(0,-1){20.00}}


\put(106.00,60.00){\line(0,-1){20.00}}


\put(73.00,60.00){\line(0,-1){20.00}}


\put(90.33,60.00){\line(0,-1){20.00}}


\put(40.00,38.00){\makebox(0,0)[cc]{$x_{1,0}$}}


\put(38.00,70.00){\makebox(0,0)[cc]{$x_2$}}


\put(38.00,60.00){\makebox(0,0)[cc]{$l_2$}}


\put(130.00,38.00){\makebox(0,0)[cc]{$x_1$}}


\put(121.00,38.00){\makebox(0,0)[cc]{$l_1$}}


\put(106.00,38.00){\makebox(0,0)[cc]{$x_{1,i+1}$}}


\put(90.33,38.00){\makebox(0,0)[cc]{$x_{1,i}$}}


\put(73.00,38.00){\makebox(0,0)[cc]{$x_{1,i-1}$}}


\put(130.00,38.00){\makebox(0,0)[cc]{$x_1$}}


\put(106.00,62.00){\makebox(0,0)[cc]{$\gamma_{i+1}$}}


\put(90.33,62.00){\makebox(0,0)[cc]{$\gamma_i$}}


\put(73.00,62.00){\makebox(0,0)[cc]{$\gamma_{i-1}$}}


\put(47.50,50.00){\makebox(0,0)[cc]{$\dotsb$}}


\put(64.00,50.00){\makebox(0,0)[cc]{${\overline\omega}_h^{\,i-1}$}}


\put(82.00,50.00){\makebox(0,0)[cc]{${\overline\omega}_h^{\,i}$}}


\put(99.00,50.00){\makebox(0,0)[cc]{${\overline\omega}_h^{\,i+1}$}}


\put(113.50,50.00){\makebox(0,0)[cc]{$\dotsb$}}


\end{picture}


\vskip -3.7cm





\centerline{Рис.~1.~Фрагмент разбиения области}





\vskip 3mm





Пусть $Y_2^{(i)}(x)$ определены при $x\in{\overline\omega}_h^{\,i}$.


В подобласти ${\overline\omega}_h^{\,i}$ определим сеточные операторы


$L_1u=-\nu h_1^{-1}u_{\ov x_1}+0,5u_{\ov


x_2x_2}+0,5h_1^{-1}a_1u


^{(-1_1)}-0,5a_2u_{\overset{\circ}{x}{}_2}$,


$L_2u=\nu h_1^{-1}u_{x_1}+0,5u_{\ov x_2x_2}-0,5h_1^{-1}a_1u


^{(+1_1)}-0,5a_2u_{\overset{\circ}{x}{}_2}$. Тогда метод декомпозиции


области для вычисления компонент скорости можно записать следующим


образом:


\begin{gather}


Y_{1t}=-\Grad_h\widehat Q+LY_2,\quad


Y_{2t}=-\Grad_hQ+L\widehat Y_2,\quad


\Div_h\widehat Y_1=0,\quad


(x,t)\in\omega_h^{i}\times\omega_\tau,\notag\\


%


\frac{\wh Y_2^{\,(i)}-0,5(Y_2^{(i-1)}+Y_2^{(i)})}{\tau}=-\Grad_hQ+\sigma


L_2(\widehat Y_2^{(i)}-Y_2^{(i)})+L_1Y_2^{(i-1)}+L_2Y_2^{(i)},\quad


(x,t)\in\gamma_{i-1}\times\omega_\tau,\\


%


\frac{\wh Y_2^{\,(i)}-0,5(Y_2^{(i+1)}+Y_2^{(i)})}{\tau}=-\Grad_hQ+\sigma


L_2(\widehat Y_2^{(i)}-Y_2^{(i)})+L_1Y_2^{(i)}+L_2Y_2^{(i+1)},\quad


(x,t)\in\gamma_{i}\times\omega_\tau.\notag


\end{gather}





Для устойчивости алгоритма (11) справедлива





\begin{thm}


Метод декомпозиции $(11)$ безусловно устойчив при $\sigma\geq 1$,


и для оценки его решения имеет место неравенство


$$


\sum_{i=0}^{N-1}(B^{(i)}(Y,Q)+b^{(i)})\leq C_5(\|\varphi_1\|_{(1)}^2+


\|\varphi_2\|_{(1)}^2+\|\Lambda\varphi_1\|^2+\|\Lambda\varphi_2\|^2),


$$


где


$B^{(i)}(Y,Q)$ --- норма $B(Y,Q)$ для подобласти $\omega_h^{(i)}$,


$b_i=\|Y_2^{i+1}-Y_2^{(i)}\|^2$ при $x\in\gamma_i$,


$C_5=\mathrm{const}>0.$


\end{thm}





Если количество подобластей конечно и не зависит от шагов сетки, то


метод


(11) сходится на гладких решениях со скоростью $O(h^2+\tau)$.


Это улучшает известные результаты по методу декомпозиции, построенные на


основе ММПН в работах [21]--[23], а также классического двухкомпонентного


метода переменных направлений, в которых оценка точности имела порядок


$h^2+h^{-1/2}\tau$. Разбиение на подобласти в методе (11) возможно


одновременно и по второй переменной. Если количество подобластей


существенно


возрастает и по порядку приближается к числу узлов пространственной


сетки,


то точность метода ухудшается, но всегда имеет порядок


$h^2+h^{-1/2}\tau$.





Анализ погрешности различных экономичных методов решения уравнений


Навье--Стокса показывает, что им присущи определенные дисбалансы,


наличие которых ведет к нарушению естественной консервативности исходной


задачи. Поэтому обычные неявные схемы, с высокой степенью адекватности


аппроксимирующие дифференциальную задачу, имеют преимущество (особенно


при


использовании ``грубых'' сеток). В связи с этим возникает вопрос о


построении


эффективных методов реализации неявных схем. Оказывается, в рамках


предложенных


выше алгоритмов можно конструировать экономичные итерационные методы,


сохраняющие все свойства исходной модели и допускающие адаптацию на


многопроцессорные ЭВМ.





Рассмотрим для уравнения (1) чисто неявную разностную схему


порядка $O(h^2+\tau)$


\begin{equation}


Y_t=-\Grad_h\widehat Q+L\widehat Y,\quad


\Div_h\widehat Y=0,


\end{equation}


с естественными начальными и граничными условиями. Для решения (12)


построим следующие итерационные процессы:


\begin{gather}


\overset{k+1}{Y}_{\!\!\!1t} =-\Grad_h\overset{k+1}{Q}{}+


L\overset{k}{Y}_{\!2},\quad


\overset{k+1}{Y}_{\!\!\!2t} =-\Grad_h\overset{k+1}{Q}{}+


L\overset{k+1}{Y}_{\!\!\!2},\quad


\Div_h\overset{k+1}{Y}_{\!\!\!2} =0,\\


%


\overset{k+1}{Y}_{\!\!\!1t} =-\sigma\Grad_h(\overset{k+1}{Q}{}-


\overset{k}{Q}{})-\Grad_h\overset{k}{Q}{}+L\overset{k}{Y}_{\!2},\notag\\


%


\overset{k+1}{Y}_{\!\!\!2t} =


\sigma L(\overset{k+1}{Y}_{\!\!\!2}-\overset{k}{Y}{}_2)-


\Grad_h \overset{k}{Q}{}+L\overset{k}{Y}_{\!2},\quad


\Div_h \overset{k+1}{Y}_{\!\!\!2}=0,


\end{gather}


где


$\overset{k+1}{Y}_{\!\!\!\alpha t} =(\overset{k+1}{Y}_{\!\!\!\alpha}


-Y)/\tau$.


Очевидно, построение итерационных методов для большего числа


компонент


разбиения, в том числе для метода декомпозиции, повторяет алгоритмы


(9)--(11).





Сходимость итерационных процессов (13), (14) доказывается аналогично


[24].


Пусть


$\overset{k+1}{\rho}_{\!\!\!\alpha} =\widehat Y-


\overset{k+1}{Y}_{\!\!\!\alpha}$, $\overset{k+1}{D}{}=-


(\widehat Q -\overset{k+1}{Q}{})$.


Справедлива





\begin{thm}


Итерационные процессы $(13)$, $(14)$ сходятся к решению неявной


разностной схемы $(12)$ со скоростью геометрической прогрессии, и для


оценки скорости сходимости имеет место неравенство


\begin{equation}


\overset{k+1}{Z}{}\leq z^{k+1}\overset{0}{Z},


\end{equation}


где


$\overset{k+1}{Z}=\|\overset{k+1}{\rho}_{\!\!1}\|_{(1)}^2+


\|\overset{k+1}{\rho}_{\!\!2}\|_{(1)}^2+


\|\Grad_h \overset{k+1}{Q}{}-


L\overset{k+1}{\rho}_{\!\!\!2}\|^2$, $z=1/(1+C_6h^2\tau^{-1})$,


$C_6>0$ --- ограниченная константа, не зависящая от $h$, $\tau$.


\end{thm}





Из оценки (15) следует, что оптимальная скорость сходимости метода


достигается при $\tau\sim h$. Это условие характерно и для обычного


метода


переменных направлений. Аналогичные итерационные процессы можно


построить


для алгоритма метода декомпозиции области (11).





Очевидно, неявная разностная схема может рассматриваться как


итерационный метод решения стационарной задачи для уравнений


Навье--Стокса


$$


\nu\sum_{\alpha=1}^2\frac{\partial^2u_k}{


\partial x_\alpha^2}+\sum_{\alpha=1}^2a_\alpha\,\frac{\partial


u_k}{\partial


x_\alpha}=-\frac{\partial P}{\partial x_k}+f_k(x),\quad k=1,2,\quad


\sum_{\alpha=1}^2\frac{\partial u_\alpha}{\partial x_\alpha}=0,\quad


x\in G,\quad u_k(x)=0,\quad x\in\Gamma.


$$


Алгоритмы (13), (14) при этом следует рассматривать как внутренние


итерационные процессы. Отметим, что увеличение размерности задачи


не создает принципиальных трудностей при построении подобных алгоритмов,


т.к. количество разбиений при использовании ММПН может быть


произвольно.





Аналогичные результаты можно получить для нелинейной задачи [3]


\begin{equation}


\frac{\partial u_k}{\partial t}-\nu\Delta u_k+


\sum_{\alpha=1}^2\frac{\partial}{\partial x_\alpha}\,(u_\alpha u_k)=


-\frac{\partial P}{\partial x_k},\quad


\sum_{\alpha=1}^2\frac{\partial u_\alpha}{\partial x_\alpha}=0


\end{equation}


с условиями (3). На основе энергетически нейтральных разностных схем


[10] для численного решения задачи (16) (в том числе и для


стационарной) строятся эффективные экономичные итерационные методы,


подобные вышеуказанным.





Как уже отмечалось выше, одно из достоинств ММПН состоит в том, что при


использовании этого подхода


в задачах декомпозиции области практически не возникает ограничений на


количество подобластей разбиения. Благодаря этому свойству на основе


ММПН


возможно построение алгоритмов пространственной декомпозиции, в которых


подобласти разбиения отождествляются с элементарными ячейками сетки.


Данный подход представляется


весьма перспективным при использовании в качестве итерационного процесса


реализации неявных разностных схем и решения стационарных задач


методом установления. Рассмотрим класс итерационных алгоритмов,


основанных на поячеечной декомпозиции задачи.





Для сшивки разностных решений в каждой из четырех вершин прямоугольной


сеточной ячейки используем аддитивное представление разностного


оператора


задачи. При этом в отличие от традиционного метода переменных


направлений


используется не покоординатное расщепление оператора, а расщепление по


диагональным направлениям. В зависимости от выбора диагонального


направления существует два варианта такого расщепления


$L=L_m+L_{m+2}$, $m=1,2$,


\begin{align*}


L_1u&=\nu(h_1^{-1}u_{x_1}+h_2^{-1}u_{x_2})-


0,5(h_1^{-1}a_1u^{(+1_1)}+h_2^{-1}a_2u^{(+1_2)}),\\


%


L_3u&=-\nu(h_1^{-1}u_{\ov x_1}+h_2^{-1}u_{\ov x_2})+


0,5(h_1^{-1}a_1u^{(-1_1)}+h_2^{-1}a_2u^{(-1_2)}),\\


%


L_2u&=-\nu(h_1^{-1}u_{\ov x_1}-h_2^{-1}u_{x_2})+


0,5(h_1^{-1}a_1u^{(-1_1)}-h_2^{-1}a_2u^{(+1_2)}),\\


%


L_4u&=\nu(h_1^{-1}u_{x_1}-h_2^{-1}u_{\ov x_2})-


0,5(h_1^{-1}a_1u^{(+1_1)}-h_2^{-1}a_2u^{(-1_2)}).


\end{align*}


Пронумеровав узлы ячеек в направлении против часовой стрелки, начиная с


нижнего левого, для совокупности вершин ячейки получим систему уравнений


для вычисления скоростей на $(k+1)$-й итерации


\begin{equation}


\frac{\overset{k+1}{Y}{\!}_{\!\!\!\!2}^{\;(i)}-


0,5\big(\widetilde Y_2^{(i+2)}+Y_2^{(i)}\big)}{\tau}=


-\Grad_hQ+\sigma L_i\big(\overset{k+1}{Y}{\!}_{\!\!\!\!2}^{\;(i)}-


\overset{k}{Y}{\!}_{\!\!2}^{\;(i)}\big)+


L_i\overset{k}{Y}{\!}_{\!\!2}^{\;(i)}+L_{i+2}


\overset{k}{\widetilde Y}{\!}_{\!\!2}^{\;(i+2)},


\end{equation}


$i=1,\dotsc,4$, $L_{i+4}\equiv L_i$, $Y_2^{(i+4)}\equiv Y_2^{(i)}$.


Здесь $\widetilde Y_2^{(i+2)}$ --- решения задач в диагонально


примыкающих


к рассматриваемой ячейке подобластях в узлах сетки, имеющих порядковые


номера $i=1,\dotsc,4$. Нетрудно заметить, что при однотипном порядке


нумерации узлов с каждым внутренним узлом будут соотнесены два


разностных


решения, принадлежащих двум диагонально примыкающим к нему ячейкам


$Y_2^{(i)},\widetilde Y_2^{(i+2)}$.





Количество уравнений в (17) совпадает с количеством неизвестных сеточных


функций на $(k+1)$-й итерации. В силу этого алгоритм замыкается на


отдельной рассматриваемой ячейке, что позволяет осуществлять расчеты в


каждой


из подобластей независимо от других. Таким образом, построен


итерационный


алгоритм явного типа реализации с максимальными возможностями для его


распараллеливания.





При использовании итерационного алгоритма (17) для отыскания


стационарного


решения задачи методом установления скорость сходимости алгоритма в


целом


(сумма итераций (17) по всем временным слоям) практически не зависит от


точности внутренних итераций (17), достигнутой на промежуточных слоях.


Это позволяет ограничиться на каждом эволюционном шаге одной итерацией,


что приводит к итерационному алгоритму сквозного счета.





В заключение коротко остановимся на оценке


вычислительных качеств многокомпонентных итерационных


алгоритмов, основанных на поячеечной декомпозиции задачи


Навье--Стокса. Наиболее существенным представляется выявление


зависимости скорости сходимости итераций от шагов пространственной


сетки,


величины итерационных параметров, а также от числа Рейнольдса.


Считается приемлемым, если для достижения заданной точности


в расчетах стационарной картины двумерных вязких течений число


итераций возрастает не быстрее чем


$O(\nu^{-1})$. Вычислительный эксперимент для модельной задачи


конвекции-диффузии показал [25], что для рассмотренного класса


методов декомпозиционного типа (17) количество итераций для


получения с заданной точностью стационарного решения


оценивается как $O(\nu^{-1/2})$. При этом наибольшая скорость


сходимости таких итерационных методов по результатам


численной проверки достигается при значении итерационных


параметров $\sigma=1$ и $\tau\sim h$ [25], что совпадает


с оценкой теоремы~5.





\begin{thebibliography}{99}





\bibitem{1}


Абрашин В.Н., Жадаева Н.Г.


{\it Разностные схемы для задач математической физики в области


произвольной формы}


// Дифференц. уравнения и их применение.- Вильнюс, 1988. -- \No\,43. --


С.\,22--30.


\bibitem{2}


Абрашин В.Н.


{\it Об одном варианте метода переменных направлений решения


многомерных задач математической физики}


// Дифференц. уравнения. -- 1990. -- Т.\,26. -- \No\,2. -- C.\,314--323.


\bibitem{3}


Абрашин В.Н., Лапко С.Л.


{\it Об одном классе разностных схем решения уравнений


Навье--Стокса.} I


// Дифференц. уравнения. -- 1992. -- Т.\,28. -- \No\,7. -- C.\,1154--1167.


\bibitem{4}


Абрашин В.Н., Лапко С.Л.


{\it Об одном классе разностных схем решения уравнений


Навье--Стокса.} II


// Дифференц. уравнения. --


1993. -- Т.\,29. -- \No\,4. -- C.\,673--688.


\bibitem{5}


Самарский А.А.


{\it Введение в теорию разностных схем.} --


М.: Наука, 1971. -- 552 с.


\bibitem{6}


Яненко Н.Н.


{\it Метод дробных шагов решения многомерных задач математической


физики.} --


Новосибирск: Наука, 1967. -- 195 с.


\bibitem{7}


Марчук Г.И.


{\it Методы расщепления.} --


М.: Наука, 1988. -- 264 с.


\bibitem{8}


Peaceman D.W., Rachford H.H.


{\it The numerical solution of parabolic and elliptic


differential equations}


// J. Soc. Industr. and Appl. Math. -- 1955. -- \No\,3. -- P.\,28--41.


\bibitem{9}


Douglas J.


{\it On the numerical integration of


$\partial^2u/\partial x^2+\partial^2u/\partial y^2=\partial u/\partial t$


by implicit methods}


// J. Soc. Industr. and Appl. Math. -- 1955. -- \No\,3. -- P.\,42-65.


\bibitem{10}


Фрязинов И.В.


{\it Консервативные разностные схемы для двумерных уравнений несжимаемой


вязкой жидкости в переменных скорость-давление} //


Препринт Ин-т прикл. матем. АН СССР. -- M., 1981. -- \No\,11. -- 28 c.


\bibitem{11}


Бахвалов Н.С., Кобельков Г.М., Чижонков Е.В.


{\it Эффективные численные методы решения уравнений Навье--Стокса}


// Числен. моделир. в аэрогидродинамике. -- M., 1986. -- С.\,37--45.


\bibitem{12}


Кобельков Г.М.


{\it Об одной разностной схеме расчета нестационарных уравнений


Навье--Стокса}


// Журн. вычисл. матем. и матем. физ. -- 1984. -- Т.\,24. -- \No\,2.


-- C.\,294--304.


\bibitem{13}


Эпштейн Б.С.


{\it Об одной схеме типа переменных направлений для задачи


Навье--Стокса}


// Вестн. Ленингр. ун-та. -- 1974. -- \No\,7. -- C.\,166--168.


\bibitem{14}


Белоцерковский О.М.


{\it Численное моделирование в механике сплошных сред.} --


М.: Наука, 1984. -- 520 с.


\bibitem{15}


Белоцерковский О.М., Гущин В.А., Щенников В.В.


{\it Метод расщепления в применении к решению задач


динамики вязкой несжимаемой жидкости}


// Журн. вычисл. матем. и матем. физ. -- 1975. -- Т.\,15. --


\No\,1. -- C.\,197--207.


\bibitem{16}


Марчук Г.И.


{\it Методы вычислительной математики.} --


М.: Наука, 1989. -- 608 с.


\bibitem{17}


Лебедев В.И.


{\it Метод композиции.} --


М.: Отд. вычисл. матем. АН СССР, 1986. -- 191 c.


\bibitem{18}


Лебедев В.И., Бахвалов Н.C., Агошков В.И. и др.


{\it Параллельные алгоритмы решения некоторых стационарных задач


математической физики.} --


М.: Отд. вычисл. матем. АН СССР, 1984. -- 142 c.


\bibitem{19}


Ладыженская О.А.


{\it Математические вопросы динамики вязкой несжимаемой жидкости.} --


М.: Наука, 1970. -- 288 с.


\bibitem{20}


Самарский А.А., Николаев Е.С.


{\it Методы решения сеточных уравнений.} --


М.: Наука, 1978. -- 592 с.


\bibitem{21}


Абрашин В.Н.


{\it Об одном методе разделения на подобласти при решении


задач математической физики}


// Дифференц. уравнения. -- 1995. -- Т.\,31. -- \No\,9. -- C.\,1525--1535.


\bibitem{22}


Самарский А.А., Вабищевич П.Н.


{\it Аддитивные схемы декомпозиции области для параболических задач}


// Дифференц. уравнения. -- 1995. -- Т.\,31. -- \No\,9. -- C.\,1563--1569.


\bibitem{23}


Лаевский Ю.М.


{\it Об одном алгоритме декомпозиции области без налегания


подобластей при решении параболических уравнений}


// Журн. вычисл. матем. и матем. физ. -- 1992. -- Т.\,32. -- \No\,11.


-- C.\,1744--1755.


\bibitem{24}


Абрашин В.Н.


{\it Об одном итерационном методе решения разностных задач


для эллиптических уравнений}


// Дифференц. уравнения. -- 1998. -- Т.\,31. -- \No\,7. -- C.\,911--920.


\bibitem{25}


Волков В.М., Лэхтиков С.Н.


{\it Многокомпонентные итерационные методы декомпозиционного типа


для двумерных стационарных задач диссипативного переноса}


// Дифференц. уравнения. -- 1997. -- Т.\,33. -- \No\,7. -- C.\,927--933.





\end{thebibliography}





\vskip 2cm





\post{\it Институт математики Национальной }{\it Поступила}


\post{\it Академии Наук Беларуси,}{06.10.1998}


\post{\it Белорусский государственный университет}{}





\label{end}


\end{document}








